\wbunit{¿Qué es una fracción?}{En esta unidad descubrirás que una fracción nos ayuda a hablar de \emph{partes iguales} de un entero: una pizza, una barra de chocolate o un camino.}

\section{Partes iguales}
Cuando partimos algo para compartirlo, lo justo es que todas las partes sean \textbf{iguales}.

\begin{wbidea}
Una fracción $\frac{a}{b}$ dice: el entero se partió en $b$ partes iguales y tomamos $a$ de ellas.
\end{wbidea}

\begin{center}
\wbfracbar{3}{4}\qquad \wbfraccircle{2}{6}\qquad \wbfracgrid{5}{2}{4}
\end{center}

\begin{wbexample}[La barra de chocolate]
Ana parte una barra en 4 partes iguales y se come 3. ¿Qué fracción se comió?

\wbfracbar[5cm]{3}{4}

Se partió en $4$ partes (denominador) y tomó $3$ (numerador): $\frac{3}{4}$.
\end{wbexample}

\begin{wbtip}
El \textbf{denominador} está \textbf{debajo} y nos dice en cuántas partes se \textbf{divide}.
\end{wbtip}

\begin{wbpractice}[Colorea y escribe]
\begin{enumerate}
\item ¿Qué fracción está coloreada? \quad \wbfraccircle[0.8cm]{1}{3} \quad \wbfracbox
\item Escribe la fracción: tres quintos. \wbblank
\item Explica con tus palabras qué significa $\frac{2}{5}$.
\wblines{2}
\end{enumerate}
\end{wbpractice}

\section{Fracciones en la recta numérica}
\wbnumberline{4}{2}{3,6}

Entre $0$ y $1$ hay cuatro saltos iguales; el punto marcado está en $\frac{3}{4}$. También $1\frac{1}{2} = \frac{6}{4}$ y $\frac{3}{4} > \frac{1}{2}$.

\begin{wbpractice}
\begin{enumerate}
\item Marca $\frac{1}{4}$ en la recta. \wbnumberline[9cm]{4}{1}{}
\item Dibuja $\frac{2}{3}$ en el recuadro. \wbbox[2.5cm]
\end{enumerate}
\end{wbpractice}

\begin{wbreview}
\begin{enumerate}
\item ¿Cuál es mayor, $\frac{1}{2}$ o $\frac{1}{4}$? \wbblank
\item Colorea $\frac{2}{5}$: \wbfracbar[6cm]{0}{5}
\end{enumerate}
\end{wbreview}
